\section{The appointed texts}
The Latin follows the 1962 Missal. The English Scripture is the historical Douay--Rheims/Challoner biblical basis, including words omitted or adapted in the chants; it is not an English translation of those adaptations. The three English orations and the Preface follow Cummiskey, Philadelphia 1861, with his abbreviated conclusions. These historical English witnesses are study texts.

\subsection{Introit}
\label{proper-introit}
\properlocus{Jer 29:11,12,14; Ps 84:2 (Hebrew 85); 1962 Missal, p.~420, no.~1749}
\begin{latinproper}
Dicit Dominus: Ego cogito cogitationes pacis, et non afflictionis: invocabitis me, et ego exaudiam vos: et reducam captivitatem vestram de cunctis locis. Ps. Benedixisti, Domine, terram tuam: avertisti captivitatem Iacob. V. Gloria Patri.
\end{latinproper}
\begin{englishwitness}{Douay--Rheims biblical basis}
29:11. For I know the thoughts that I think towards you, saith the
Lord, thoughts of peace, and not of affliction, to give you an end and
patience.

29:12. And you shall call upon me, and you shall go: and you shall pray
to me, and I will hear you.

29:14. And I will be found by you, saith the Lord: and I will bring
back your captivity, and I will gather you out of all nations, and from
all the places to which I have driven you out, saith the Lord: and I
will bring you back from the place to which I caused you to be carried
away captive.

Ps 84:2. Lord, thou hast blessed thy land: thou hast turned away the
captivity of Jacob.
\end{englishwitness}

\subsection{Collect}
\label{proper-collect}
\properlocus{1962 Missal, p.~420, no.~1750}
\begin{latinproper}
Familiam tuam, quaesumus, Domine, continua pietate custodi: ut, quae in sola spe gratiae caelestis innititur, tua semper protectione muniatur. Per Dominum nostrum Iesum Christum, Filium tuum: Qui tecum vivit et regnat in unitate.
\end{latinproper}
\begin{englishwitness}{Cummiskey 1861}
Preserve, we beseech thee, O Lord, thy family by thy constant mercy, that we, who confide solely in the support of thy heavenly grace, may be always defended by thy protection. Thro'.
\end{englishwitness}

\subsection{Epistle}
\label{proper-epistle}
\properlocus{Col 3:12--17; 1962 Missal, p.~420, no.~1751}
\begin{latinproper}
Fratres: Induite vos sicut electi Dei, sancti et dilecti, viscera misericordiae, benignitatem, humilitatem, modestiam, patientiam: supportantes invicem, et donantes vobismetipsis, si quis adversus aliquem habet querelam: sicut et Dominus donavit vobis, ita et vos. Super omnia autem haec, caritatem habete, quod est vinculum perfectionis: et pax Christi exsultet in cordibus vestris, in qua et vocati estis in uno corpore: et grati estote. Verbum Christi habitet in vobis abundanter, in omni sapientia, docentes, et commonentes vosmetipsos, psalmis, hymnis, et canticis spiritualibus, in gratia cantantes in cordibus vestris Deo. Omne quodcumque facitis in verbo aut in opere, omnia in nomine Domini Iesu Christi, gratias agentes Deo et Patri per Iesum Christum Dominum nostrum.
\end{latinproper}
\begin{englishwitness}{Douay--Rheims biblical basis}
3:12. Put ye on therefore, as the elect of God, holy and beloved, the
bowels of mercy, benignity, humility, modesty, patience:

3:13. Bearing with one another and forgiving one another, if any have a
complaint against another. Even as the Lord hath forgiven you, so do
you also.

3:14. But above all these things have charity, which is the bond of
perfection.

3:15. And let the peace of Christ rejoice in your hearts, wherein also
you are called in one body: and be ye thankful.

3:16. Let the word of Christ dwell in you abundantly: in all wisdom,
teaching and admonishing one another in psalms, hymns and spiritual
canticles, singing in grace in your hearts to God.

3:17. All whatsoever you do in word or in work, do all in the name of
the Lord Jesus Christ, giving thanks to God and the Father by him.
\end{englishwitness}

\subsection{Gradual}
\label{proper-gradual}
\properlocus{Ps 43:8--9 (Hebrew 44); 1962 Missal, p.~420, no.~1752}
\begin{latinproper}
Liberasti nos, Domine, ex affligentibus nos: et eos, qui nos oderunt, confudisti. V. In Deo laudabimur tota die, et in nomine tuo confitebimur in saecula.
\end{latinproper}
\begin{englishwitness}{Douay--Rheims biblical basis}
43:8. But thou hast saved us from them that afflict us: and hast put
them to shame that hate us.

43:9. In God shall we glory all the day long: and in thy name we will
give praise for ever.
\end{englishwitness}

\subsection{Alleluia}
\label{proper-alleluia}
\properlocus{Ps 129:1--2 (Hebrew 130); 1962 Missal, p.~420, no.~1753}
\begin{latinproper}
Alleluia, alleluia. V. De profundis clamavi ad te, Domine: Domine, exaudi orationem meam. Alleluia.
\end{latinproper}
\begin{englishwitness}{Douay--Rheims biblical basis}
129:1. Out of the depths I have cried to thee, O Lord:

129:2. Lord, hear my voice. Let thy ears be attentive to the voice of
my supplication.
\end{englishwitness}

\subsection{Gospel}
\label{proper-gospel}
\properlocus{Matt 13:24--30; 1962 Missal, p.~421, no.~1754}
\begin{latinproper}
In illo tempore: Dixit Iesus turbis parabolam hanc: Simile factum est regnum caelorum homini, qui seminavit bonum semen in agro suo. Cum autem dormirent homines, venit inimicus eius, et superseminavit zizania in medio tritici, et abiit. Cum autem crevisset herba, et fructum fecisset, tunc apparuerunt et zizania. Accedentes autem servi patrisfamilias, dixerunt ei: Domine, nonne bonum semen seminasti in agro tuo? Unde ergo habet zizania? Et ait illis: Inimicus homo hoc fecit. Servi autem dixerunt ei: Vis, imus, et colligimus ea? Et ait: Non: ne forte colligentes zizania, eradicetis simul cum eis et triticum. Sinite utraque crescere usque ad messem, et in tempore messis dicam messoribus: Colligite primum zizania, et alligate ea in fasciculos ad comburendum, triticum autem congregate in horreum meum.
\end{latinproper}
\begin{englishwitness}{Douay--Rheims biblical basis}
13:24. Another parable he proposed to them, saying: The kingdom of
heaven is likened to a man that sowed good seed in his field.

13:25. But while men were asleep, his enemy came and oversowed cockle
among the wheat and went his way.

13:26. And when the blade was sprung up, and had brought forth fruit,
then appeared also the cockle.

13:27. And the servants of the good man of the house coming said to
him. Sir, didst thou not sow good seed in thy field? Whence then hath
it cockle?

13:28. And he said to them: An enemy hath done this. And the servants
said to him: Wilt thou that we go and gather it up?

13:29. And he said: No, lest perhaps gathering up the cockle, you root
up the wheat also together with it.

13:30. Suffer both to grow until the harvest, and in the time of the
harvest I will say to the reapers: Gather up first the cockle, and bind
it into bundles to burn, but the wheat gather ye into my barn.
\end{englishwitness}

\subsection{Offertory}
\label{proper-offertory}
\properlocus{Ps 129:1--2 (Hebrew 130); 1962 Missal, p.~421, no.~1755}
\begin{latinproper}
De profundis clamavi ad te, Domine: Domine, exaudi orationem meam: de profundis clamavi ad te, Domine.
\end{latinproper}
\begin{englishwitness}{Douay--Rheims biblical basis}
129:1. Out of the depths I have cried to thee, O Lord:

129:2. Lord, hear my voice. Let thy ears be attentive to the voice of
my supplication.
\end{englishwitness}

\subsection{Secret}
\label{proper-secret}
\properlocus{1962 Missal, p.~421, no.~1756}
\begin{latinproper}
Hostias tibi, Domine, placationis offerimus: ut et delicta nostra miseratus absolvas, et nutantia corda tu dirigas. Per Dominum.
\end{latinproper}
\begin{englishwitness}{Cummiskey 1861}
We offer thee, O Lord, this sacrifice of propitiation: that thou wouldst mercifully forgive us our sins, and guide our tottering hearts. Thro'.
\end{englishwitness}

\subsection{Preface of the Trinity}
\label{proper-preface}
\properlocus{Appointment: 1962 Missal, p.~421; text p.~293, no.~1082}
\begin{latinproper}
V. Per omnia saecula saeculorum. R. Amen.

V. Dominus vobiscum. R. Et cum spiritu tuo.

V. Sursum corda. R. Habemus ad Dominum.

V. Gratias agamus Domino Deo nostro. R. Dignum et iustum est.

Vere dignum et iustum est, aequum et salutare, nos tibi semper et ubique gratias agere: Domine, sancte Pater, omnipotens aeterne Deus: Qui cum unigenito Filio tuo, et Spiritu Sancto, unus es Deus, unus es Dominus: non in unius singularitate personae, sed in unius Trinitate substantiae. Quod enim de tua gloria, revelante te, credimus, hoc de Filio tuo, hoc de Spiritu Sancto, sine differentia discretionis sentimus. Ut in confessione verae sempiternaeque Deitatis, et in personis proprietas, et in essentia unitas, et in maiestate adoretur aequalitas. Quam laudant Angeli atque Archangeli, Cherubim quoque ac Seraphim: qui non cessant clamare cotidie, una voce dicentes:

Sanctus, Sanctus, Sanctus Dominus Deus Sabaoth. Pleni sunt caeli et terra gloria tua. Hosanna in excelsis. Benedictus qui venit in nomine Domini. Hosanna in excelsis.
\end{latinproper}
\begin{englishwitness}{Cummiskey 1861}
P. World without end. R. Amen.

P. The Lord be with you. R. And with thy spirit. P. Lift up your hearts. R. We have lifted them up to the Lord. P. Let us give thanks to the Lord our God. R. It is meet and just.

IT is truly meet and just, right and available to salvation, that we should always, and in all places, give thanks to thee, O holy Lord, Father Almighty, eternal God. Who together with thy only begotten Son and the Holy Ghost, art one God, and one Lord not in a singularity of one Person, but in a Trinity of one substance. For what we believe of thy glory, as thou hast revealed, the same we believe of thy Son, and of the Holy Ghost, without any difference or distinction. So that in the confession of the true and eternal Deity, we adore a distinction in the Persons, an unity in the essence, and an equality in the Majesty. Whom the angels and archangels, the cherubim also and seraphim praise: and cease not daily to cry out with one voice saying, Holy, holy, holy, Lord God of Sabaoth. Heaven and earth are full of thy glory. Hosanna in the highest. Blessed is he that cometh in the name of the Lord, Hosanna in the highest.
\end{englishwitness}

\subsection{Communion}
\label{proper-communion}
\properlocus{Mark 11:24; 1962 Missal, p.~421, no.~1757}
\begin{latinproper}
Amen dico vobis, quidquid orantes petitis, credite quia accipietis, et fiet vobis.
\end{latinproper}
\begin{englishwitness}{Douay--Rheims biblical basis}
11:24. Therefore I say unto you, all things, whatsoever you ask when ye
pray, believe that you shall receive: and they shall come unto you.
\end{englishwitness}

\subsection{Postcommunion}
\label{proper-postcommunion}
\properlocus{1962 Missal, p.~421, no.~1758}
\begin{latinproper}
Quaesumus, omnipotens Deus: ut illius salutaris capiamus effectum, cuius per haec mysteria pignus accepimus. Per Dominum.
\end{latinproper}
\begin{englishwitness}{Cummiskey 1861}
We beseech thee, O Almighty God, that we may one day receive the effects of that salvation, of which we have received the pledge in these mysteries. Thro'.
\end{englishwitness}
